\section*{Chapter \ref{ch:m1:exchanges-brokers-venues} --- Exchanges, Brokers and Venues}

\begin{solution}{exo:m1:exchanges-brokers-venues:1}
Of \$5\,411 million: energy 40.3\%; cash equities and equity options 8.6\%;
data and connectivity 19.1\%; listings 9.1\%.
\end{solution}

\begin{solution}{exo:m1:exchanges-brokers-venues:2}
$2 \times 0.70/280\,000 = 5\times10^{-6}$: \qty{0.05}{\bp}. A stock trade on
the same value would pay far more in spread alone.
\end{solution}

\begin{solution}{exo:m1:exchanges-brokers-venues:3}
(a) \omterm{def:m1:exchanges-brokers-venues:broker}{Broker}'s algorithm or desk; the \omterm{def:m1:exchanges-brokers-venues:broker}{broker}'s checks. (b) \omterm{def:m1:exchanges-brokers-venues:dma}{Direct market
access}; the \omterm{def:m1:exchanges-brokers-venues:broker}{broker}'s checks, in the \omterm{def:m1:exchanges-brokers-venues:broker}{broker}'s systems. (c) \omterm{def:m1:exchanges-brokers-venues:dma}{Sponsored access};
still the \omterm{def:m1:exchanges-brokers-venues:broker}{broker}'s checks, under its exclusive control, although they run
beside the client's machines. (d) Own membership; the firm's own checks, for
which it answers to its regulator and to its clearing firm.
\end{solution}

\begin{solution}{exo:m1:exchanges-brokers-venues:4}
\$180\,000; \$65\,000; \$49\,000; \$150\,000. \omterm{def:m1:exchanges-brokers-venues:dma}{Sponsored access} saves
\$16\,000 a month over DMA. Membership wins above 312.5 million shares a
month, five times this firm's volume.
\end{solution}

\begin{solution}{exo:m1:exchanges-brokers-venues:5}
DMA to sponsored: $20\,000/(0.0010-0.0006) = 50$ million shares a month
(was 33). Sponsored to membership: $125\,000/0.0006 = 208$ million (was
313). The range in which \omterm{def:m1:exchanges-brokers-venues:dma}{sponsored access} is best shrinks from both ends: the
provider loses its smallest clients to DMA and pushes its largest to become
members --- the ones who paid it most.
\end{solution}

\begin{solution}{exo:m1:exchanges-brokers-venues:6}
$x^\star = 40/(504+50) = 7.2\%$; at 10\%, $0.10 \times 554 - 40 = \$15.4$
million, against \$40.6 million before. A cut of one hundredth of a cent
removes three fifths of the profit.
\end{solution}

\begin{solution}{exo:m1:exchanges-brokers-venues:7}
Only \omterm{def:m1:exchanges-brokers-venues:dma}{sponsored access} (50 microseconds) and membership (40) meet 1
millisecond. At 5 million shares a month \omterm{def:m1:exchanges-brokers-venues:dma}{sponsored access} costs \$27\,000
against \$150\,000. Unconstrained, the choice would be DMA at \$10\,000: the
latency requirement costs \$17\,000 a month, 0.34 cent a share --- more than
the \omterm{def:m1:exchanges-brokers-venues:broker}{broker}'s algorithm charged for everything.
\end{solution}

\begin{solution}{exo:m1:exchanges-brokers-venues:8}
The plan compares fees and leaves out the \emph{price} obtained and the
probability of obtaining it. A taker's cost per share is the fee plus the
half-spread available on the venue, and an order that cannot be filled there
costs the delay and the move of the price meanwhile. The comparison is
$\text{fee}_A + s_A/2$ against $\text{fee}_B + s_B/2$, weighted by fill
probability. The established venue has the orders, so its spread is a cent
where the newcomer's book is empty or five cents wide: a saving of 0.25 cent
in fees against 2 cents more in spread. Without liquidity the cheap venue is
the expensive one, and a low taker fee does nothing to attract the resting
orders that would change this.
\end{solution}

\begin{solution}{pb:m1:exchanges-brokers-venues:1}
\textbf{1.} $0.28 - 0.25 = 0.03$ cent.
\textbf{2.} $0.0003/45 = \qty{0.067}{\bp}$.
\textbf{3.} $8\times10^9 \times 252 \times 0.0003 = \$604.8$ million; \$6.05
million per point.
\textbf{4.} $6.05 + 0.60 = \$6.65$ million per point.
\textbf{5.} $38/664.8 = 5.72\%$.
\textbf{6.} 457 million shares, \$20.6 billion a day.
\textbf{7.} $12\% \times \tfrac13 \times \tfrac12 = 2\%$.
\textbf{8.} No: $2 \times 6.65 - 38 = -\$24.7$ million a year.
\textbf{9.} $\nu = -0.04$ cent: the venue pays to match. At 4\%:
$0.04 \times 8\times10^9 \times 252 \times (-0.0004) = -\$32.3$ million.
\textbf{10.} $-32.3 + 2.4 - 38 = -\$67.9$ million.
\textbf{11.} $6 \times 6.65 - 38 = \$1.9$ million: thirty-six such years to
repay one year of subsidy. At 6\% the venue is alive, not valuable.
\textbf{12.} $38/(403.2 + 60) = 8.2\%$.
\textbf{13.} As owners they pay themselves part of the fees they would pay
anyway, and --- the real reason --- a credible alternative disciplines the
incumbent's fees on the other 94\% of their volume. A venue that loses
\$25 million but takes 0.02 cent off incumbent fees on the consortium's flow
can pay for itself.
\textbf{14.} It gives up the right to list, its own share of consolidated
data revenue and the regulatory status that lets it set binding rules; it
saves the cost and delay of \omterm{def:m1:exchanges-brokers-venues:exchange}{exchange} registration and self-regulation.
\textbf{15.} It attracts resting orders from slower institutions and \omterm{def:m1:what-a-trading-firm-does:market-maker}{market
makers} who fear being picked off, since the delay lets them reprice first; it
repels the fastest takers, whose edge is exactly that race.
\textbf{16.} Each trader's benefit from a venue grows with the orders already
there, so share feeds on itself: a few venues, each dominant in a niche
(a fee model, an order type, a client group), survive; the rest stay below
$x^\star$ and close or are bought.
\textbf{17.} If data revenue is allocated partly by time spent quoting at
the best price, orders that sit at the best price but are cancelled before
they can trade earn revenue for the venue without providing liquidity; the
venue may pay for them through rebate tiers.
\textbf{18.} That each sponsor is a member in good standing with a clearing
arrangement, has pre-trade controls on the sponsored flow under its exclusive
control, and can cut the client off; the venue also applies its own
order-level limits and kill switch per identifier.
\textbf{19.} \textbf{\$20.6 billion a day.}
\textbf{20.} Regulatory and commercial: the matching technology is a
commodity, while the licence, the data rights, and above all the order flow
committed by its owners decide whether it lives.
\end{solution}

\begin{solution}{iq:m1:exchanges-brokers-venues:1}
Listing fees, transaction fees net of rebates, clearing fees and income on
collateral where it owns the clearing house, market-data licences, and access
(ports, colocation, connectivity). For a high-frequency firm the transaction
line is often close to zero or negative after rebates; what it pays for is
data and access, the fixed and rising part.
\iqlookfor{the five lines, and awareness that the fast firm's bill is data
and connectivity.}
\end{solution}

\begin{solution}{iq:m1:exchanges-brokers-venues:2}
DMA orders travel through the \omterm{def:m1:exchanges-brokers-venues:broker}{broker}'s systems; sponsored-access orders go
straight from the client to the \omterm{def:m1:exchanges-brokers-venues:exchange}{exchange} under the \omterm{def:m1:exchanges-brokers-venues:broker}{broker}'s identifier. In
both the \omterm{def:m1:exchanges-brokers-venues:broker}{broker} is responsible to the \omterm{def:m1:exchanges-brokers-venues:exchange}{exchange} and the regulator: since the
market access rule it must have pre-trade controls on that flow under its
direct and exclusive control, so the erroneous orders are its failure of
control as well as the client's bug --- and its money if the client cannot
pay.
\iqlookfor{responsibility follows the membership identifier.}
\end{solution}

\begin{solution}{iq:m1:exchanges-brokers-venues:3}
Per order: maximum quantity and notional; price collar against a reference
price; instrument permitted and tradable; short-sale and regulatory flags.
Aggregate: position and notional limits per instrument and overall, open
order exposure, order and message rates, loss limit, duplicate-order and
self-match protection, plus a kill switch. The budget is microseconds or
less, so checks are in-process, branch-light arithmetic on state kept
incrementally --- no locks, no allocation, no database or network call in
the path --- and their limits are loaded at start and changed only through an
audited side channel.
\iqlookfor{aggregate as well as per-order checks, and design consequences of
the budget.}
\end{solution}

\begin{solution}{iq:m1:exchanges-brokers-venues:4}
A share is fungible across venues and the law forces venues to compete for
it, so matching fees fell toward cost. A futures contract is the \omterm{def:m1:exchanges-brokers-venues:exchange}{exchange}'s
own product and clears in its clearing house: open interest cannot move, so
liquidity cannot either. The futures \omterm{def:m1:exchanges-brokers-venues:exchange}{exchange} is a monopolist in each
successful contract and is priced like one.
\iqlookfor{fungibility and clearing as the mechanism, not ``futures are
bigger''.}
\end{solution}

\begin{solution}{iq:m1:exchanges-brokers-venues:5}
Compare expected profit per unit of time, not rebate per fill: fill rate
$\times$ (half-spread $+$ rebate $-$ \omterm{def:m1:what-a-trading-firm-does:adverse-selection}{adverse selection}). A fifth of the
fills needs five times the \omterm{prop:m1:what-a-trading-firm-does:capture}{net capture} to compete, and a rebate 0.12 cent
higher does not provide it --- unless the fills on the new venue are much
less toxic, which is a measurable question (markouts by venue). Quote on
both if capital and risk limits allow; the rebate may also count toward a
volume tier.
\iqlookfor{markouts by venue, and thinking in rates, not per-fill
economics.}
\end{solution}

\begin{solution}{iq:m1:exchanges-brokers-venues:6}
Fair: the products are offered to anyone at a published price; those who buy
them are the ones who keep quotes tight and consistent across venues, and
investors get better prices as a result; speed is a cost of doing that
business like any other. Unfair: the \omterm{def:m1:exchanges-brokers-venues:exchange}{exchange} sells, to some, an advantage
over its other customers that exists only because it also publishes a slower
version to them; the fees are set by a monopolist over its own data; and the
resulting race consumes resources in shaving microseconds with no gain to
end investors, while the slow feed is what best-execution obligations are
measured against.
\iqlookfor{both sides argued properly; the candidate's own view matters less
than seeing the structure.}
\end{solution}
